\section{Penalty-Based Baselines}
\label{sec:penalty}

\subsection{Classical Penalty QUBO}

The standard approach to enforcing the cardinality constraint $\sum_i x_i = k$ in quantum optimization is to add a quadratic penalty to the objective:
\begin{equation}
\label{eq:penalty_qubo}
E_{\mathrm{pen}}(x) = C_G(x) + \lambda \left( \sum_{i=1}^n x_i - k \right)^2, \quad x \in \{0,1\}^n,
\end{equation}
where $\lambda > 0$ is a penalty coefficient. For sufficiently large $\lambda$, the ground state of $E_{\mathrm{pen}}$ coincides with the constrained optimum. In practice, $\lambda$ must be tuned: too small allows unbalanced solutions; too large creates a rugged landscape with large coefficient ratios that challenge quantum hardware.

Expanding the penalty term gives a QUBO with all-to-all couplings:
\begin{equation}
E_{\mathrm{pen}}(x) = \sum_{(i,j)\in E} w_{ij} (x_i + x_j - 2x_i x_j) + \lambda \left( \sum_i x_i^2 + 2\sum_{i<j} x_i x_j - 2k\sum_i x_i + k^2 \right).
\end{equation}
Since $x_i^2 = x_i$ for binary variables, this yields linear terms $h_i = \sum_{j:(i,j)\in E} w_{ij} - \lambda(2k-1)$ and quadratic couplings $J_{ij} = -2w_{ij}\mathbb{1}_{(i,j)\in E} + 2\lambda$ for $i<j$. The all-to-all connectivity and coefficient ratio $\max|J_{ij}|/\min|h_i| \sim \lambda/w$ can be problematic for hardware with limited dynamic range.

\subsection{Photonic Penalty: ObliQ}

The ObliQ algorithm~\cite{wang2023bosonic} implements a photonic version of penalty-based optimization. It uses $n$ modes with $N \ge k$ photons (typically $N=n$), and a threshold detection scheme where each mode's photon number is binarized: $x_i = \mathbb{1}[n_i > 0]$. The penalty is implemented by allowing the interferometer to explore the full binary space $\{0,1\}^n$, with the penalty term evaluated classically from the thresholded outcomes.

Key differences from our fixed-$k$ approach:
\begin{itemize}
\item \textbf{Photon number:} ObliQ uses $N \ge k$ photons (often $N=n$), not fixed $k$. The output photon number is not conserved in the thresholded space.
\item \textbf{Search space:} Full $\{0,1\}^n$ (dimension $2^n$) vs. our collision-free subspace (dimension $\binom{n}{k}$).
\item \textbf{Objective:} Penalized expectation $\mathbb{E}[C_G(x) + \lambda(\sum x_i - k)^2]$ vs. our $P_{\mathrm{success}}$.
\item \textbf{Postselection:} ObliQ discards no outcomes (all thresholded bitstrings are kept); we postselect on collision-free events.
\end{itemize}

\subsection{Classical Product Bernoulli Baseline}

For fair comparison at equal classical evaluation budget, we implement a classical product Bernoulli sampler with the same penalized objective. The model has $n$ independent logits $\theta_i$ defining probabilities $q_i = \sigma(\theta_i)$ for each bit. The probability of bitstring $x$ is $\prod_i q_i^{x_i}(1-q_i)^{1-x_i}$. The expected penalized energy is minimized via gradient-based optimization (COBYLA or Adam). This baseline has the same number of trainable parameters ($n$ logits) as our Clements mesh ($n(n-1)/2 + n$ parameters) and operates on the full binary space.

\subsection{Comparison Metrics}

For fair comparison at equal \textbf{physical sampling cost}, we evaluate:
\begin{itemize}
\item \textbf{Physical success probability:} $P_{\mathrm{success}} = \sum_{x: \sum x_i=k, C_G(x)=C^*} p(x)$.
\item \textbf{Expected shots:} $N_{\mathrm{shots}} = 1/P_{\mathrm{success}}$.
\item \textbf{Equal budget:} Same number of classical objective evaluations (e.g., 200) and same number of physical shots per evaluation.
\end{itemize}

Our variational fixed-$k$ approach uses the collision-free subspace as a \textbf{physical constraint} (enforced by photon number conservation), while penalty methods use a \textbf{soft constraint} (enforced by energy penalty). The former restricts the search space but guarantees balanced outputs upon collision-free detection; the latter explores a larger space but requires postselection or large $\lambda$ to enforce balance.
